% Integration fragment; statement labels are provisional and should be renumbered by Part.
% AI-produced by GPT-6 Astra (OpenAI, ChatGPT Pro mode). No priority claim.

\subsection{What the two frame signs do and do not identify}

\paragraph{A.2-L1 (Refuted; corrected below): the frames are central.}
Write $\mathcal H=\ker(\lambda:\widehat{\mathrm{GT}}_0\to
\widehat{\mathbf Z}^{\times})$, and use precisely the signed lifts and
cusp correction of LM7\_PUSH E.65.  The proved maps in that statement are
\[
 T:\mathcal H\longrightarrow
 C_{\mathrm{SL}_2(\mathbf F_3)}(R_0)
 =\langle-I\rangle\times\langle R_0\rangle,
 \qquad
 W:\mathcal H\longrightarrow
 C_{\mathrm{SL}_2(\mathbf F_3)}(S)=\langle S\rangle,
\]
with $S=\left(\begin{smallmatrix}0&1\\-1&0\end{smallmatrix}\right)$.
These are centralizers, not the scalar centre.  In particular the assertion
that both corrected frames are scalar modulo twelve is false even on
$G_{\mathbf Q}\cap\mathcal H$.

Here is an actual Galois counterexample, specified by a nonempty compact
coset of the absolute Galois group.  Put $r_0=(3/4)^{1/4}>0$.
The polynomial $4X^4-3$ is Eisenstein at $3$, so
$[\mathbf Q(r_0):\mathbf Q]=4$.  This real field is not normal over
$\mathbf Q$: it does not contain the root $ir_0$ of that irreducible
polynomial.  Every finite subextension of $\mathbf Q^{\rm ab}/\mathbf Q$
is normal, and consequently $r_0\notin\mathbf Q^{\rm ab}$.
On the other hand $r_0^2=\sqrt3/2$ belongs to $\mathbf Q^{\rm ab}$.
Thus $\mathbf Q^{\rm ab}(r_0)/\mathbf Q^{\rm ab}$ is quadratic.
The restriction map
\[
 G_{\mathbf Q^{\rm ab}}\longrightarrow
 \operatorname{Gal}(\mathbf Q^{\rm ab}(r_0)/\mathbf Q^{\rm ab})
\]
is a continuous surjection.  Take any element $\sigma$ in the inverse
image of its nonidentity element.  This inverse image is nonempty and
compact, and $\sigma(r_0)=-r_0$ and $\chi(\sigma)=1$.

The signed Galois comparison in the referee's repaired proof of
LM7\_PUSH E.65, using the two specific models
$E_3:y^2=x^3+1$ and $E_4:y^2=x^3-x$, gives
\[
 \epsilon_3(T_\sigma)=\epsilon_2(W_\sigma)=1.
\]
Therefore $W_\sigma\bmod3$ is $S$ or $-S$.  It is not scalar, even
projectively.  Indeed, for $U=\left(\begin{smallmatrix}1&1\\0&1\end{smallmatrix}\right)$,
\[
 SU=\begin{pmatrix}0&1\\2&2\end{pmatrix},\qquad
 US=\begin{pmatrix}2&1\\2&0\end{pmatrix}\quad\text{over }\mathbf F_3,
\]
and these matrices are neither equal nor negatives of each other.
Reduction modulo three of a central scalar class modulo twelve would be
scalar.  This proves the claimed counterexample.  It refutes the literal
central-frame assertion, not the parity compression that follows.

\paragraph{A.2-L2 (Proved): CRT compression and its precise ambiguity.}
Let $a(\Phi)=\epsilon_3(T_\Phi)$ and
$b(\Phi)=\epsilon_2(W_\Phi)$ for $\Phi\in\mathcal H$.
There is a unique continuous homomorphism
\[
 u:\mathcal H\longrightarrow\mu_2(\mathbf Z/12)
 =\{1,5,7,11\}
\]
whose reductions modulo three and four are respectively $(-1)^a$ and
$(-1)^b$.  Explicitly,
\[
\begin{array}{c|c|c|c}
(a,b)&u&e\text{ with }e^2=e,\;2e=1-u&\delta=a+b\\\hline
(0,0)&1&0&0\\
(1,0)&5&4&1\\
(0,1)&7&9&1\\
(1,1)&11&1&0
\end{array}
\]
Continuity and multiplicativity follow from LM7\_PUSH E.65 and CRT.
The idempotents of $\mathbf Z/12\simeq\mathbf F_3\times\mathbf Z/4$
are the four pairs with both coordinates zero or one; the table follows.
The expression $(1-u)/2$ is not ordinary division in $\mathbf Z/12$.
Each displayed equation $2e=1-u$ has two solutions in that ring, and
exactly one is idempotent.  This condition is essential.

The element of R2 has $u=5$.  Any independently established lift with
parities $(0,1)$ has $u=7$; the existence of that lift must be proved
separately and is not a consequence of this table.  On
$G_{\mathbf Q}\cap\mathcal H$ the values lie in $\{1,11\}$, and the
preceding counterexample shows that $11$ really occurs.

The ring $\mathbf F_3\times\mathbf Z/4$ combines the two \emph{signs}.
It does not combine the two actual frames as one scalar matrix.  Their
natural adelic centralizer algebras are different quadratic algebras,
the Eisenstein and Gaussian algebras.  Passing from their units to the
two indicated characters loses data.

For an arbitrary $\Phi\in\widehat{\mathrm{GT}}_0$, choose
$\sigma\in G_{\mathbf Q}$ with $\chi(\sigma)=\lambda(\Phi)$.
The well-defined output is
\[
 \overline u(\Phi)=
 [u(\Phi\sigma^{-1})]\ \in\
 \{1,5,7,11\}/\{1,11\}.
\]
If $\sigma$ is changed, the two arguments differ by an element of
$G_{\mathbf Q}\cap\mathcal H$ after conjugation, so the corresponding
sign-pairs differ by the diagonal sign.  The parity sum and its
conjugacy invariance, proved in LM7\_PUSH E.65, show that this formula
is independent of the choice, continuous, and multiplicative.
It corresponds to $\Delta(\Phi)$ under the unique isomorphism of its
target with $\mathbf F_2$.  The full lift $u(\Phi)$ has not been
canonically defined for general multiplier by this construction.

\paragraph{A.2-L3 (Proved): the theta identity and an induced action.}
For $r\in\mathbf Z/12$ define
\[
 \theta_{6,r}(\tau,z)=
 \sum_{n\equiv r\;(12)}q^{n^2/24}e^{2\pi inz},
 \qquad q=e^{2\pi i\tau},\quad\operatorname{Im}\tau>0.
\]
Let $P_u$ permute the labels by $r\mapsto ur$.
All four units $u\in\{1,5,7,11\}$ satisfy $u^2=1\pmod{24}$;
thus they preserve the finite quadratic module
$A=\mathbf Z/12$, $q_A(r)=r^2/24\pmod{\mathbf Z}$.
In particular $P_u$ commutes with its diagonal $T$ matrix and its
finite Fourier $S$ matrix.  Euler's pentagonal identity gives
\[
 \eta(\tau)=q^{1/24}\prod_{n\ge1}(1-q^n)
 =\sum_{n\ge1}\chi_{12}(n)q^{n^2/24}
 =\theta_{6,1}(\tau,0)-\theta_{6,5}(\tau,0),
\]
where $\chi_{12}$ is zero on nonunits, positive at $1,11$, and
negative at $5,7$.  For example, multiplying Euler's exponent
$k(3k-1)/2$ by $24$ and adding $1$ gives $(6k-1)^2$;
sorting the positive representatives modulo twelve gives the stated
Dirichlet character.  Absolute convergence permits the rearrangement.
Also $\theta_{6,-r}(\tau,0)=\theta_{6,r}(\tau,0)$, so direct
substitution proves
\[
 P_u\eta=\chi_{12}(u)\eta,
 \qquad
 \chi_{12}(u(\Phi))=(-1)^{\delta(\Phi)}.
\]
At $z=0$ these are twelve labelled functions with linear dependencies;
on the formal twelve-dimensional label space the corresponding
eigenvector is $e_1+e_{11}-e_5-e_7$.

It follows that $\Phi\mapsto P_{u(\Phi)}$ is a continuous
representation of $\mathcal H$ on this label space and that its eta
line carries $(-1)^\delta$.  In this \emph{defined} representation
the R2 element acts by $P_5$ and gives eta eigenvalue $-1$; a lift
with parities $(0,1)$ acts by $P_7$ and has the same eta eigenvalue.
No comparison with an independently defined modular-functor action
has been proved.  This construction packages the known character;
it does not derive its restriction to $\widehat{\mathrm{GT}}$.

\paragraph{A.2-L4 (Proved): a pointed category kills the $f$-coordinate.}
We can make the issue completely explicit.  Work over
$K=\mathbf Q(\zeta_{24})$.  Use representatives $0,\ldots,11$
for $A=\mathbf Z/12$, and put
\[
 [a+b]=a+b-12\kappa(a,b),\qquad
 \kappa(a,b)=\left\lfloor\frac{a+b}{12}\right\rfloor.
\]
Consider finite-dimensional $A$-graded $K$-vector spaces with the
usual graded tensor product, associator and braiding on homogeneous
simple objects given by
\[
 \omega(a,b,c)=(-1)^{a\kappa(b,c)},\qquad
 c(a,b)=\zeta_{24}^{ab}.
\]
The pentagon follows from the carry identity
\[
 \kappa(b,c)+\kappa([b+c],d)
 =\kappa(c,d)+\kappa(b,[c+d]).
\]
More explicitly the exponent difference between its two sides is
$12\kappa(a,b)\kappa(c,d)$ modulo two.  For the first hexagon,
$c(a,[b+c])/(c(a,b)c(a,c))=(-1)^{a\kappa(b,c)}$,
and the associator corrections give the same value; the factors
containing $\kappa(a,c)=\kappa(c,a)$ cancel.  For the second,
$c([a+b],c)/(c(a,c)c(b,c))=(-1)^{c\kappa(a,b)}$,
and the same cancellation proves the equality.  This is a braided
pointed category with trace quadratic form
$c(a,a)=\zeta_{24}^{a^2}$ and nondegenerate double braiding
$\zeta_{12}^{ab}$.

For every pair of its double-braiding operators, evaluation of a
word $f\in\widehat F_2'$ is the identity.  Indeed the operators
are simultaneously diagonal on the homogeneous components of every
tensor product, with eigenvalues in the finite abelian group
$\mu_{12}$.  The unique continuous evaluation homomorphism from
$\widehat F_2$ to that finite diagonal group therefore kills
$\widehat F_2'$.  This statement holds on direct sums as well as
on individual simple objects.

Drinfeld's formula for the change of associator multiplies it by
such an evaluation of $f$ (\S4, equations (4.5), (4.10), and (4.12)
in the English edition).  Hence this correction is exactly one.
His braiding change, for $m=(\lambda-1)/2$, multiplies $c(a,b)$ by
the $m$th power of the double braiding, giving
\[
 c_\lambda(a,b)=c(a,b)^{\lambda}
 =\zeta_{24}^{\lambda ab},\qquad \omega_\lambda=\omega.
\]
These expressions make sense for every profinite unit $\lambda$
by reduction modulo twenty-four.  Since $\lambda$ is odd,
$\omega^\lambda=\omega$.  They also follow by applying the
field automorphism $\zeta_{24}\mapsto\zeta_{24}^{\lambda}$ to
the original coherence equations.  Therefore they satisfy both
hexagons and the pentagon for \emph{every} $f\in\widehat F_2'$,
including every gentle pair.  At $\lambda=1$ the resulting category
and its chosen structural maps are literally unchanged.

With the ribbon convention in which the trace quadratic form is
the twist, its modular data are
\[
 T_{ab}=\delta_{ab}\zeta_{24}^{a^2},\qquad
 S_{ab}=12^{-1/2}\zeta_{12}^{ab}.
\]
The opposite Fourier convention replaces the exponent by its
negative and makes no difference to the conclusion.  Cyclotomic
conjugation transforms $T$ by entrywise $\lambda$th powers.  If the
normalizing square root $\sqrt{12}=2\sqrt3$ is conjugated too,
then
\[
 \sigma_\lambda(S)_{ab}
 =\chi_{12}(\lambda)\,12^{-1/2}\zeta_{12}^{\lambda ab}
 =\chi_{12}(\lambda)S_{\lambda a,b}.
\]
The sign is obtained directly from
$\sqrt3=\zeta_{12}+\zeta_{12}^{-1}$.
In particular the pointed-category modular data depend on
$\lambda$, not on $f$.

\paragraph{A.2-L5 (Refuted as a test; Open as a comparison).}
The assertion that the R2 element must fail the categorical pentagon
in the indicated pointed category is false.  Its multiplier is one,
so the previous paragraph proves that its proposed modification is
the identity modification and satisfies that pentagon.
This does \emph{not} prove that R2 satisfies the group pentagon.
The evaluation functor used here kills every $f$-word in the
commutator subgroup, and thus also kills the information that this
test was supposed to detect.

The distinction in Drinfeld's result is the quantifier: imposing
the group pentagon makes the construction coherent on all suitable
braided categories, naturally in braided functors.  A single
pointed category is a nonfaithful test of that condition.  The
unproved implication needed for the proposed route is an
identification, on the eta line or on the full label space, of
the actual frame representation with the independent modular
functor operation of a GT pair.  In our notation it would identify
$P_{u(\Phi)}$ with that operation after the specified cusp and
Galois normalizations.  For multiplier one the independent operation
is the identity whereas $P_{u(\Phi)}$ can be nonidentity.  Proving
that identification on all of $\widehat{\mathrm{GT}}\cap\mathcal H$
would already imply $\Delta=0$.  It is not supplied by the
category's pentagon or by the equality of modular-group
representations.  Consequently this route gives no decision of
LM7\_PUSH E.3.

\paragraph{A.2-L6 (Proved for homomorphically lifted frames): abelianized phases cannot detect the bit.}
The projection to $B^{\rm ab}\simeq\mathbf Z/6$ cannot recover the
R2 parity from its pure frame.  Indeed the image of the level-two
free subgroup $F$ in $B^{\rm ab}$ has order three: with
$\tau=sr$ and $x=\tau^{-2}$, $y=sxs$, both generators map to
$-2\tau$.  Thus $F'F^3$ maps to zero.  The identity frame and
R2's frame in $F'F^3$ have the same abelianized image but different
parities.  The following additional argument establishes the
metaplectic version with an explicit, choice-independent convention.

In fact every multiplier-one positive $r$-frame belongs to $F'$.
The identities $\Phi(r)=f^{-1}r=ArA^{-1}$ and $f\in F'$
give $(I-R)[A]=0$ in $F^{\rm ab}=\widehat{\mathbf Z}^2$, where
conjugation by $r$ sends $(x,y,z)$ to $(y,z,x)$ and hence
\[
 R=\begin{pmatrix}0&-1\\1&-1\end{pmatrix},\qquad
 I-R=\begin{pmatrix}1&1\\-1&2\end{pmatrix},\qquad\det(I-R)=3.
\]
Multiplication by three is injective on $\widehat{\mathbf Z}$,
so the adjugate identity proves $[A]=0$. For the R2 element also
$\beta\in R_{\rm nil}(F'F^3)\subseteq(F'F^3)'\subseteq F'$,
because its retained order-two conjugator is one. Its shear is
therefore trivial, so its two corrected frames and their inverses
belong to $F'$.

Let $p:E\twoheadrightarrow\widehat B$ be any profinite central
extension, in particular the completion of
$\operatorname{Mp}_2(\mathbf Z)\to\operatorname{PSL}_2(\mathbf Z)$.
Choose lifts in $E$ of the two free generators of $F$. Freeness
gives a continuous homomorphism $\iota:F\to E$ lifting the
inclusion into $\widehat B$. For two such homomorphic lifts,
their pointwise ratio is a homomorphism from $F$ to the central
kernel, and therefore vanishes on $F'$. Thus $\iota$ restricted
to $F'$ is independent of these choices. In particular both
R2 frames have a well-defined lifted value, and their images
in $E^{\rm ab}$ are zero because $\iota$ is a homomorphism.
The identity has the same two abelianized images and bit zero,
whereas R2 has bit one. No function of these abelianized
metaplectic frame values can equal $\Delta$.

One may realize the usual metaplectic extension explicitly as
$B_3/\langle z_3^4\rangle$, mapping to
$B_3/\langle z_3\rangle=B$, with central kernel of order four;
its abelianization is $\mathbf Z/24$. The argument above does
not require this presentation. It also shows why arbitrary
pointwise lifts, which can be changed by a central element,
must not be confused with the homomorphic lift just specified.
The scalar eta multiplier and the label permutation $P_u$
are different operations even after this ambiguity is removed.

\paragraph{A.2-L7 (Known and Proved): the scalar eta phase and its central table.}
For $M=\left(\begin{smallmatrix}a&b\\c&d\end{smallmatrix}\right)
\in\mathrm{SL}_2(\mathbf Z)$ with $c>0$, the classical formula is
\[
 \eta(M\tau)=
 \exp\!\left(\pi i\left(\frac{a+d}{12c}-s(d,c)-\frac14\right)\right)
 (c\tau+d)^{1/2}\eta(\tau),
\]
with the usual holomorphic square-root branch and
$s(d,c)=\sum_{k=1}^{c-1}((k/c))((dk/c))$.
This is the Dedekind-sum scalar multiplier; the values at the
generators are
\[
 \eta(\tau+1)=\zeta_{24}\eta(\tau),\qquad
 \eta(-1/\tau)=(-i\tau)^{1/2}\eta(\tau).
\]
See Atiyah (1987), equations (1.3)--(1.4), for the logarithmic
version, and Theorem 5.60 for its relation to signature defects,
eta invariants, and the Meyer cocycle. Those comparisons concern
the scalar automorphy factor on the modular group.

If a metaplectic element is written $(M,\varphi)$ with
$\varphi(\tau)^2=c\tau+d$, then
$\eta(M\tau)/(\varphi(\tau)\eta(\tau))$ is its scalar
character $\nu_\eta$. The cocycle law for $\varphi$ proves
multiplicativity directly. Its square descends to a character
$\nu_{\eta^2}:\mathrm{SL}_2(\mathbf Z)\to\mu_{12}$,
with $T\mapsto\zeta_{12}$ and
$S=\left(\begin{smallmatrix}0&-1\\1&0\end{smallmatrix}\right)
\mapsto\zeta_{12}^{-3}$.
The attached exact enumeration constructs a homomorphism
\[
 e:\mathrm{SL}_2(\mathbf Z/12)\longrightarrow\mathbf Z/12,
 \qquad e(T)=1,\quad e(S)=-3.
\]
It enumerates all $1152$ determinant-one matrices and verifies all
$4608$ right-generator edges, proving well-definedness and
multiplicativity. Equality on $S,T$ proves
$\nu_{\eta^2}(M)=\zeta_{12}^{e(M\bmod12)}$. Thus on the scalar
centre the exact table is
\[
\begin{array}{c|rrrr}
 u&1&5&7&11\\\hline
 \nu_{\eta^2}(uI)&1&1&-1&-1\\
 \chi_{12}(u)&1&-1&-1&1
\end{array}
\]
and the first row is $\chi_{-4}(u)$. This proves the proposed
eta-squared central-character calculation, including the distinction
between $5$ and $7$. It does not identify the actual frames with
these central matrices: A.2-L1 refutes that identification. On the
homomorphically lifted R2 frames every scalar character is trivial
by A.2-L6, while the separately defined label permutation has eta
eigenvalue $-1$. The remaining genus-zero/genus-one comparison is
therefore still the one in A.2-L5.

\subsection{The AA condition and finite characters}

\paragraph{A.3-L1 (Known and Proved): AA subgroups lie in character kernels.}
In Hoshi--Mochizuki--Tsujimura (2025), ``AA'' means
\emph{algebraically ample}, with Definition 5.12 including the
COF, RGC, and ArBC prerequisites.  The proof of Theorem 5.17(ii)
in the author version dated 14 March 2025, printed page 78,
states that every closed AA subgroup of
$G=\widehat{\mathrm{GT}}$ is contained in a $G$-conjugate of
$G_{\mathbf Q}$; the maximal such subgroups are exactly those
conjugates.  The theorem is not an equality $G=G_{\mathbf Q}$.

Let $\eta:G\to C$ be any continuous homomorphism to an abelian
topological group which is trivial on $G_{\mathbf Q}$.  Then
\[
 \eta(g\sigma g^{-1})
 =\eta(g)\eta(\sigma)\eta(g)^{-1}=1.
\]
Thus every closed AA subgroup lies in $\ker\eta$.  This proves
the requested assertion for $\Delta|_G$ and for any additional
characters of A.3(h) if those characters are actually constructed.
No assertion about the existence of the latter is required.

The marked assertion in R10 also follows.  By
Lochak--Schneps (1997), Proposition 4(ii),
$C_G(c)=\langle c\rangle$ for $c=(-1,1)$.
Suppose $gG_{\mathbf Q}g^{-1}$ contains $c$.  All involutions of
$G_{\mathbf Q}$ are conjugate to complex conjugation, so there
is $h\in G_{\mathbf Q}$ with
$g^{-1}cg=hch^{-1}$.  Hence $gh\in C_G(c)=\langle c\rangle
\subseteq G_{\mathbf Q}$, so $g\in G_{\mathbf Q}$ and
$gG_{\mathbf Q}g^{-1}=G_{\mathbf Q}$.  Consequently
$G_{\mathbf Q}$ is the unique maximal closed AA subgroup
containing the marked $c$.  This proof uses no classification of
involutions of $G$ outside its Galois subgroup.

\paragraph{P-L1 (Proved): the precise abelianization question.}
For either $G=\widehat{\mathrm{GT}}$ or
$G=\widehat{\mathrm{GT}}_0$ let
$G^{\rm ab}=G/\overline{[G,G]}$ and
$K_G=\ker(G^{\rm ab}\xrightarrow{\lambda}\widehat{\mathbf Z}^{\times})$.
Kronecker--Weber identifies
$G_{\mathbf Q}^{\rm ab}$ with $\widehat{\mathbf Z}^{\times}$
via $\chi$.  The inclusion $G_{\mathbf Q}\hookrightarrow G$
therefore gives a continuous section of the displayed multiplier
map.  Thus
\[
 G^{\rm ab}\simeq\widehat{\mathbf Z}^{\times}\times K_G
\]
as profinite abelian groups, with the Galois abelianization as the
first factor.  Every continuous finite abelian character of
$G_{\mathbf Q}$ extends to $G$ by composition with $\lambda$.
Characters of $G$ trivial on $G_{\mathbf Q}$ are exactly the
continuous characters of $K_G$.  For finite-order complex
characters this is its Pontryagin dual.

Lochak--Schneps, \emph{Open problems in Grothendieck--Teichm\"uller
theory}, Question 1.2, asks whether $K_{\widehat{\mathrm{GT}}}=1$.
Equivalently, whether every finite-order character of
$\widehat{\mathrm{GT}}$ factors through $\lambda$: characters
separate points of a profinite abelian group, since every nonidentity
element survives in a finite abelian quotient and in a character
of that quotient.  R4 proves that
$K_{\widehat{\mathrm{GT}}_0}$ has a quotient $C_2$.
This last fact neither computes that kernel nor decides whether
the restriction of this quotient to $\widehat{\mathrm{GT}}$
is nonzero.  We found no classification of the finite characters
or resolution of Question 1.2 in the primary sources searched.

